\documentclass[11pt]{article}

\usepackage[a4paper,margin=1in]{geometry}
\usepackage{amsmath,amssymb,amsthm}
\usepackage{booktabs}
\usepackage[hidelinks]{hyperref}
\usepackage{microtype}

\newtheorem{theorem}{Theorem}
\newtheorem{lemma}{Lemma}
\newtheorem{proposition}{Proposition}
\newtheorem{remark}{Remark}
\newcommand{\dd}{\,\mathrm{d}}
\newcommand{\one}{\mathbf{1}}
\DeclareMathOperator{\sinc}{sinc}

\title{The Erd\H{o}s Minimum-Overlap Constant Exceeds 0.38055925}
\author{Ruturaj R Raval\\
\small Independent Researcher\\
\small ORCID: \href{https://orcid.org/0000-0003-4930-8981}
{0000-0003-4930-8981}}
\date{4 September 2026}

\begin{document}
\maketitle

\begin{abstract}
We prove the computer-assisted lower bound \(c_E>0.38055925\) for
Erd\H{o}s' minimum-overlap constant. Price's publicly released 172-bin Arb
certificate is retained on 170 noncentral mean bins, while its two central
bins are replaced by a new even dual certificate using one second-moment
inequality, 107 pointwise Fourier inequalities, and a Parseval inequality
truncated at order 100. The positive-part integral of the new certificate is
bounded by two independently implemented directed-arithmetic checkers, one
using Arb through python-flint and one using MPFI, MPFR, and GMP. A fresh run
of Price's original Arb checker verifies every retained noncentral bin at the
same target. The optimization used to discover the multipliers is outside the
trusted proof. Exact certificate data, verification code, hashes, and run
records are archived with the paper. The exact value of \(c_E\) remains open.
\end{abstract}

\section{Introduction}

Let
\[
  \{1,\ldots,2N\}=A\sqcup B,\qquad |A|=|B|=N,
\]
and write
\[
  R(A,B)=\max_{d\in\mathbb Z}
  |\{(a,b)\in A\times B:b-a=d\}|.
\]
The Erd\H{o}s minimum-overlap constant is
\[
  c_E=\liminf_{N\to\infty}
  \frac1N\min_{\substack{A\sqcup B=\{1,\ldots,2N\}\\|A|=|B|=N}}
  R(A,B).
\]
Erd\H{o}s introduced the problem in 1955 \cite{Erdos1955}. Moser obtained the
classical lower bound \cite{Moser1959}, and White later introduced a stronger
Fourier-analytic relaxation \cite{White2023}. Recent computer-assisted work
has strengthened the lower bound through finite dual certificates
\cite{Chung2026,Price2026,Deng2026}. The exact value remains unknown.

The present result is a finite, computer-assisted lower-bound proof. It
builds on Price's mean-bin framework and noncentral certificates
\cite{Price2026}, and on the center-bin Fourier formulation used by Deng
\cite{Deng2026}. The project-original content is a stronger center
certificate, two independently implemented directed-arithmetic checkers for
that certificate, and the combined global audit. The optimization that found
the multipliers is not part of the trusted proof.

\section{Continuous formulation}

Let $\mathcal H$ be the set of measurable functions
$h:[0,2]\to[0,1]$ with $\int_0^2 h=1$, and put $g=1-h$. Extend both functions
by zero outside $[0,2]$. Define
\[
  p_h(t)=\int_{\mathbb R}h(x)g(x+t)\dd x,\qquad -2\le t\le2.
\]
Then $p_h\ge0$ and $\int_{-2}^{2}p_h=1$. Set
\[
  c_{\mathrm{cont}}=\inf_{h\in\mathcal H}\|p_h\|_\infty.
\]

\begin{lemma}\label{lem:transfer}
One has $c_E\ge c_{\mathrm{cont}}$.
\end{lemma}

\begin{proof}
For a balanced partition of $\{1,\ldots,2N\}$, replace each integer by its
interval of length $1/N$ in $[0,2]$. The resulting overlap density is
\[
  p(t)=\frac1N\sum_{a\in A,\,b\in B}
  (1-|Nt-(b-a)|)_+.
\]
For fixed $Nt$, at most two adjacent integer differences contribute, and
their triangular weights sum to at most one. Hence
$\|p\|_\infty\le R(A,B)/N$. Taking the minimum and then the liminf proves the
claim.
\end{proof}

It therefore suffices to lower-bound $\|p_h\|_\infty$ uniformly over
$h\in\mathcal H$.

\section{Analytic constraints}

Fix $h\in\mathcal H$, write $p=p_h$, and define
\[
  \mu=\int_{-2}^{2}t\,p(t)\dd t.
\]

\begin{lemma}\label{lem:meanrange}
Every admissible overlap density satisfies \(\mu\in[-1,1]\).
\end{lemma}

\begin{proof}
Let
\[
  m_h=\int_0^2 xh(x)\dd x,\qquad
  m_g=\int_0^2 xg(x)\dd x.
\]
Since \(h+g=1\), one has \(m_h+m_g=2\). Also
\[
  \int_0^1(1-h)=\int_1^2h=:\delta.
\]
Using \(x\le1\) on \([0,1]\) and \(x\ge1\) on \([1,2]\),
\[
  m_h-\frac12
  =\int_0^1x(h-1)\dd x+\int_1^2xh\dd x
  \ge-\delta+\delta=0.
\]
Applying the same argument to \(g\) gives \(m_g\ge1/2\), hence
\(m_h\le3/2\). Finally, the change of variables \(y=x+t\) gives
\[
  \mu=m_g-m_h=2-2m_h,
\]
which lies in \([-1,1]\).
\end{proof}

\subsection{Second moment}

\begin{lemma}\label{lem:second}
For every admissible $p$,
\[
  \int_{-2}^{2}t^2p(t)\dd t=\frac23+\frac{\mu^2}{2}.
\]
Consequently, if $|\mu|\le1/320$, then
\[
  \int t^2p(t)\dd t\le
  B_2:=\frac23+\frac{1}{2\cdot320^2}.
\]
\end{lemma}

\begin{proof}
Let $X$ and $Y$ be independent with densities $h$ and $g$. Since $h+g=1$,
\[
  \mathbb EX+\mathbb EY=2,\qquad
  \mathbb EX^2+\mathbb EY^2=\frac83.
\]
Also $\mu=\mathbb E(Y-X)$, so
$\mathbb EX=1-\mu/2$ and $\mathbb EY=1+\mu/2$. Therefore
\[
  \mathbb E(Y-X)^2
  =\frac83-2\mathbb EX\,\mathbb EY
  =\frac23+\frac{\mu^2}{2}.
\]
\end{proof}

\subsection{Fourier rows}

For real $\xi$, define
\[
  H(\xi)=\int_0^2h(x)e^{i\xi(x-1)}\dd x,\qquad
  G(\xi)=\int_0^2g(x)e^{i\xi(x-1)}\dd x,
\]
and
\[
  C_\xi=\int_{-2}^{2}p(t)\cos(\xi t)\dd t,\qquad
  \sinc\xi=\begin{cases}\sin\xi/\xi,&\xi\ne0,\\1,&\xi=0.\end{cases}
\]
Then $H(\xi)+G(\xi)=2\sinc\xi$ and the Fourier transform of $p$ is
$\overline{H(\xi)}G(\xi)$.

\begin{lemma}\label{lem:cosine}
For every real $\xi$,
\[
  C_\xi\le\sinc^2\xi.
\]
\end{lemma}

\begin{proof}
Write $H-G=a+ib$ and $\sigma=\sinc\xi$. Since $H+G=2\sigma$,
\[
  \overline HG
  =\sigma^2-\frac{a^2+b^2}{4}-i\sigma b.
\]
Taking real parts gives the claim.
\end{proof}

\begin{lemma}[Parseval row]\label{lem:parseval}
For every positive integer $M$,
\[
  -\sum_{n=1}^{M}C_{n\pi}\le\frac12.
\]
\end{lemma}

\begin{proof}
At $\xi=n\pi$, $n\ne0$, one has $G(\xi)=-H(\xi)$ and therefore
$C_{n\pi}=-|H(n\pi)|^2$. Parseval on $[-1,1]$ gives
\[
  \sum_{n\in\mathbb Z}|H(n\pi)|^2
  =2\int_0^2h(x)^2\dd x\le2.
\]
Since $H(0)=1$ and the positive and negative frequencies have equal squared
modulus,
\[
  \sum_{n=1}^{\infty}|H(n\pi)|^2\le\frac12.
\]
Every finite truncation obeys the same upper bound.
\end{proof}

\section{Positive-part dual bound}

\begin{lemma}\label{lem:dual}
Suppose $p\ge0$, $\int p=1$, and
\[
  \int p(t)\phi_j(t)\dd t\le B_j,\qquad j=1,\ldots,r.
\]
For nonnegative multipliers $\lambda_j$, define
\[
  q(t)=1+\sum_{j=1}^{r}\lambda_j(B_j-\phi_j(t)).
\]
Then
\[
  \|p\|_\infty\ge
  \left(\int_{-2}^{2}\max(q(t),0)\dd t\right)^{-1}.
\]
\end{lemma}

\begin{proof}
The constraints and nonnegative multipliers give $\int pq\ge1$. Since $p\ge0$,
\[
  1\le\int pq\le\int p\max(q,0)
  \le\|p\|_\infty\int\max(q,0).
\]
\end{proof}

\section{The center certificate}

For the two bins
\[
  I_-=[-1/320,0],\qquad I_+=[0,1/320],
\]
the retained certificate defines
\begin{align*}
q(t)=1
&+\lambda_2(B_2-t^2)\\
&+\sum_{\xi\in\Xi}
  \lambda_\xi\bigl(\sinc^2\xi-\cos(\xi t)\bigr)\\
&+\lambda_P\left(\frac12+\sum_{n=1}^{100}\cos(n\pi t)\right).
\end{align*}
The exact decimal values of the nonnegative multipliers and the 107
frequencies in $\Xi$ are stored in
\texttt{certificates/center-038055925.tsv}. The certificate is even, so the
same function covers both center bins.

Every decimal frequency and multiplier in the file is interpreted as an
exact rational number. For a cosine row, however, the right-hand side is the
exact transcendental value $\sinc^2\xi$, not a decimal approximation. Each
verifier independently encloses this value with directed arithmetic.

Both parsers require the declared target, mean interval, row counts, and
Parseval order; reject malformed or duplicate rows; and verify that every
multiplier is nonnegative. Each checker recomputes every $\sinc^2\xi$ with
outward-rounded arithmetic. Neither checker trusts the discovery objective,
reported roots, or optimizer termination status.

\subsection{Rigorous positive-part integration}

On an interval $J=[a,b]$ with center $c$ and radius $r$, the verifier encloses
$q(c)$, $q'(c)$, and $q''(c)$. It also computes a global directed upper bound
$M_3$ for $\sup_{[-2,2]}|q'''|$. More precisely, it evaluates the sum
\[
  S_3=
  \sum_{\xi\in\Xi}\lambda_\xi\xi^3
  +\lambda_P\sum_{n=1}^{100}(n\pi)^3
\]
with directed arithmetic and takes \(M_3\) to be its outward-rounded upper
endpoint. The triangle inequality therefore gives
\[
  \sup_{t\in[-2,2]}|q'''(t)|\le S_3\le M_3.
\]
Taylor's theorem gives
\[
  q(x)\in q(c)+[-E_J,E_J],
\]
where
\[
  E_J=|q'(c)|r+\frac{|q''(c)|r^2}{2}+\frac{M_3r^3}{6}.
\]
All quantities in this formula are interval enclosures. On a
certified-positive interval, the checker evaluates the antiderivative
\begin{align*}
Q(t)=t
&+\lambda_2\left(B_2t-\frac{t^3}{3}\right)\\
&+\sum_{\xi\in\Xi}\lambda_\xi
  \left(\sinc^2\xi\,t-\frac{\sin(\xi t)}{\xi}\right)\\
&+\lambda_P\left(\frac t2+
  \sum_{n=1}^{100}\frac{\sin(n\pi t)}{n\pi}\right)
\end{align*}
with outward rounding.

If the resulting upper endpoint is nonpositive, the interval contributes
zero. If the lower endpoint is nonnegative, the elementary antiderivative of
$q$ is evaluated with directed arithmetic. Otherwise the interval is
bisected. An unresolved terminal interval is charged by its width times the
positive part of the certified upper endpoint. The calculation is performed
on $[0,2]$ and doubled by evenness.

\begin{proposition}\label{prop:center}
The canonical certificate satisfies
\[
  \int_{-2}^{2}\max(q(t),0)\dd t
  \le 2.627711172296609115765958739219.
\]
\end{proposition}

\begin{proof}
The Python-Arb checker, at 256-bit precision with 4096 initial cells and
maximum bisection depth 16, returns the displayed outward-rounded upper
endpoint. It classifies 4805 cells as negative and 5903 as positive, and
charges 120 terminal cells by upper rectangles. The complete command,
certificate hash, directed endpoints, and cell counts are retained in the
machine-readable evidence record.
\end{proof}

In particular,
\[
  2.627711172296609115765958739219
  <\frac1{0.38055925}
  =2.627711716375308181\ldots.
\]

\subsection{Independent MPFI implementation}

A second verifier is implemented in C with MPFI, MPFR, and GMP. It has its own
strict certificate parser, exact decimal-to-rational conversion, interval
evaluation, adaptive traversal, antiderivative, accounting, and endpoint
extraction. At 256-bit precision it gives
\[
  \int_{-2}^{2}\max(q(t),0)\dd t
  \le2.627711172296609115765958701269.
\]
This independently establishes the same strict target margin. The two
implementations share the certificate bytes and mathematical formulas, but
not parser code, arithmetic backend, traversal code, or result extraction.
This is implementation independence within the project, not an external
third-party review.

\section{Noncentral bins}

Price's public certificate partitions the full mean interval \([-1,1]\) into
172 bins \cite{Price2026}. The only bins replaced here are 85 and 86, namely
\(I_-\) and \(I_+\). The retained source state is commit
\texttt{6bc610e40083ef61a40966dfb5d38612cabc4c5b}. The retained SHA-256 identities
are:
\begin{center}
\begin{minipage}{0.94\linewidth}
\small
certificate:\quad
\texttt{a2c723d375073bd124c04a690200ca9bf8f598eac506656e1ac9a883e6f4f217}

\smallskip
checker:\quad
\texttt{a825b28f0dd8b8db3d8f225bcd8af48374c013f2f59de97bdb3d8d8cd62a49f11e}
\end{minipage}
\end{center}

The project reran Price's original Arb checker at target \(0.38055925\) on
bins \(0\) through \(84\) and \(87\) through \(171\). All 170 checks passed.
The source decimals left a gap from \(0.025\) to
\(0.025000000000000133\). Before verification, bin 94 was widened outward by
replacing its lower endpoint with \(0.025\). The checker then recomputed every
mean-dependent right-hand side from the widened endpoints, so the repair
enlarges the certified domain and omits no mean value. The largest certified
denominator occurs at bin 77 and its reflected partner 94:
\[
  D_{\mathrm{noncentral}}
  \le2.627538530873375790019512034721.
\]
Thus
\[
  D_{\mathrm{noncentral}}
  <\frac1{0.38055925}.
\]
Because Price's repository had no declared license at the pinned commit, its
source and certificate are not redistributed here. Only provenance, hashes,
the exact invocation, mathematical facts, and the aggregate replay record are
retained. This rerun authenticates the cited dependency but is not a
clean-room implementation of Price's noncentral checker. The configuration
audit verifies the settings used here, including positive right-hand-side
inflation \(10^{-14}\) and 160-bit Arb arithmetic; it does not certify every
parameter combination accepted by the general-purpose source program.

\section{Global result}

\begin{theorem}\label{thm:main}
The Erd\H{o}s minimum-overlap constant satisfies
\[
  \boxed{c_E>0.38055925}.
\]
\end{theorem}

\begin{proof}
Let $h\in\mathcal H$ and let $\mu$ be the mean of $p_h$. If
$|\mu|\le1/320$, Lemmas \ref{lem:second}, \ref{lem:cosine},
\ref{lem:parseval}, and \ref{lem:dual}, together with
Proposition \ref{prop:center}, give
\[
  \|p_h\|_\infty>0.38055925.
\]
Otherwise $\mu$ lies in one of Price's 170 retained noncentral bins, whose
certified denominator is also strictly below $1/0.38055925$. Hence
the finite maximum of the certified denominators over all 172 bins is strictly
below $1/0.38055925$. This gives the uniform bound
$c_{\mathrm{cont}}>0.38055925$. Lemmas \ref{lem:meanrange} and
\ref{lem:transfer} complete the proof.
\end{proof}

\section{Verification and scope}

The trusted center proof consists of the canonical certificate, exact decimal
semantics, the analytic lemmas above, and two directed positive-part
implementations. The trusted noncentral dependency consists of Price's pinned
certificate and checker identities together with a fresh run of the original
checker on every retained bin. The optimization, exploratory floating-point
root searches, and untraced solver output are outside the trusted path.

This theorem improves the lower bound but does not determine $c_E$. A gap
remains to the best reported upper constructions. The noncentral part depends
on Price's public certificate, the retained multipliers are not claimed to be
optimal, and independent external mathematical review remains desirable. No
claim is made that rerunning Price's checker constitutes an independently
implemented noncentral proof.

\section{Data and code availability}

The exact project release preceding this paper is archived at
\href{https://doi.org/10.5281/zenodo.22308924}
{doi:10.5281/zenodo.22308924}. It contains the center certificate, both center
checkers, machine-readable directed endpoints, source hashes, the aggregate
170-bin replay record, tests, and reproduction commands. The present arXiv
ancillary bundle additionally contains the complete project-generated
per-bin CSV, JSON, and log. Price's source and certificate are not
redistributed because the pinned repository state had no declared license.
Their repository, commit, paths, hashes, and exact replay command are recorded
so that a reviewer can obtain them from the original source.

\section*{Acknowledgements}

The author thanks Liam Price for publicly releasing the mean-bin certificate
and checker on which the noncentral part depends. The author also acknowledges
the Fourier-analytic framework of Ethan White, the central-bin refinement of
Haowen Deng, and the publicly released certificate artifacts of Stephen
Chung, Wenyu Du, and William J. Wesley.

\begin{thebibliography}{9}

\bibitem{Erdos1955}
P. Erd\H{o}s,
\emph{Some remarks on number theory},
Riveon Lematematika 9 (1955), 45--48.

\bibitem{Moser1959}
L. Moser,
\emph{On the minimum overlap problem of Erd\H{o}s},
Acta Arithmetica 5 (1959), 117--119.

\bibitem{White2023}
E. P. White,
\emph{A new bound for Erd\H{o}s' minimum overlap problem},
Acta Arithmetica 208 (2023), no.~3, 235--255,
doi:10.4064/aa220728-7-6; arXiv:2201.05704.

\bibitem{Chung2026}
S. Chung, W. Du, and W. J. Wesley,
\emph{Autonomous mathematical discovery in an open-world multi-agent
environment},
arXiv:2608.23691, 2026.

\bibitem{Price2026}
L. Price,
\emph{A 0.38055470 lower bound for Erd\H{o}s Problem 36},
GitHub repository,
\url{https://github.com/Leeham06972452/erdos-36-lower-bound},
commit \texttt{6bc610e}, 2026.

\bibitem{Deng2026}
H. Deng,
\emph{A sharper certified lower bound for Erd\H{o}s' minimum-overlap problem},
version 1.0.1, Zenodo, doi:10.5281/zenodo.22279895, 2026.

\end{thebibliography}

\end{document}
